\chapter*{Notation}\label{ch:notation}

En este libro utilizamos la notación de una variedad de disciplinas matemáticas,
que se resumen a continuación. Se recomienda al lector que consulte esta lista como referencia a medida de que vaya leyendo los diversos capítulos del libro.
Hemos tratado de mantener un empleo consistente de estas notaciones a lo largo
de los capítulos.

Como un primer ejemplo, para un número natural \(n\) el conjunto 
\(\{1, 2, \dots, n\}\) se denota por \([n]\).

\section*{Notación de álgebra lineal}
\begin{itemize}
    \item Los escalares (ya sean deterministas o aleatorios) no están en negrita
        p.ej., \(x\), \(X\).
    \item Los vectores (ya sea deterministas o aleatorios) están en negritas
        minúsculas, p.ej., \(\vx\), \(\vpi\). Los elementos de vectores se 
        denotan \(x_{i}\) o \(\pi_{i}\), alternativamente \((\vx)_{i}\).
        Los términos genéricos (que pueden ser escalares, vectores, matrices o
        tensores de orden superior) también están en negritas minúsculas.
    \item Las matrices (ya sean deterministas o aleatorios) están en negritas 
        mayúsculas, p.ej., \(\vX, \vPi\). Los elementos de matrices se denotan
        \(X_{ij}\) o \(\Pi_{ij}\), alternativamente \((\vX)_{ij}\).
        Las columnas de matrices son vectores en negritas y minúsculas, p.ej.,
        \(\vx_{i}\) es la \(i\th\) columna de \(\vX\), a no ser que se defina 
        de otra manera.
    \item La transpuesta de una matriz es \(\top\), p.ej.,
        \(\vA^{\top}\). La adjunta es \(\adj\), p.ej., \(\vA\adj\).
        La pseudo-inversa es \(\dagger\), p.ej., \(\vA^{\dagger}\).
    \item El rango de una matriz es \(\rank(\vA)\), la traza es
        \(\tr(\vA)\), el determinante es \(\det(\vA)\) y su
        determinante logarítmico es \(\logdet(\vA)\).
    \item La multiplicación elemento por elemento de matrices es \(\vA \hada
        \vB\). La elevación al cuadrado elemento por elemento es 
        \(\vA^{\hada 2}\), etc. De manera similar, el producto Kronecker de
        matrices es \(\vA \kron \vB\). El producto Kronecker iterado
        es \(\vA^{\kron 2}\), etc.
    \item Para una función \(f \colon \R \to \R\), su aplicación elemento por
        elemento a una matriz \(\vX\) es \(f[\vX]\), esto es, con
        paréntesis cuadrados. Para una función \(f \colon \R^{d} \to
        \R^{k}\) y una matriz \(\vX \in \R^{d \times n}\), podemos
        difundir \(f\) a fin de aplicarla a cada columna sin una
        notación especial, esto es, \(f(\vX) = [f(\vx_{1}), \dots, f(\vx_{n})]
        \in \R^{k \times n}\).
    \item La esfera unitaria (euclidiana) es \(\Sphere^{n - 1} \subseteq
        \R^{n}\). la bola unitaria (euclidiana) es \(\Ball^{n} \subseteq \R^{n}\).
    \item El conjunto de matricas ortogonales es \(\O(m, n) \subseteq
        \R^{m \times n}\) o \(\O(n) = \O(n, n)\).
    \item Las matrices simétricas son \(\Sym(n)\), las matrices PSD simétricas
        son \(\PSD(n)\), las matrices PD simétricas son \(\PD(n)\).
    \item Los valores propios de una matriz simétrica \(\vS \in \Sym(n)\) son
        todos reales según el teorema espectral; se denotan
        \(\lambda_{i}(\vS)\) y se ordenan de modo que \(\lambda_{1}(\vS)
        \geq \cdots \geq \lambda_{n}(\vS)\).
    \item Los valores singulares de \(\vA \in \R^{m \times n}\) con
        \(\rank(\vA) = r\) se denotan  \(\sigma_{i}(\vA)\) y se
        ordenan de modo que \(\sigma_{1}(\vA) \geq \cdots \geq
        \sigma_{r}(\vA) > 0\). Por convención, definimos \(\sigma_{r +
        1}(\vA) = \sigma_{r + 2}(\vA) = \cdots = 0\).
    \item La proyección sobre un conjunto \(\cK\) se denota \(\proj_{\cK}\).
    \item Las normas \(\ell^{p}\) de vectores se denotan por
        \(\norm{\vx}_{p}\), \(p \in [0, \infty]\).
    \item La norma euclidiana de operadores sobre matrices es is \(\norm{\vA} =
        \sigma_{1}(\vA)\).
    \item La norma Frobenius de matrices es \(\norm{\vA}_{F} =
        \tr(\vA^{\top}\vA)\).
    \item El producto interno euclidiano de vectores es \(\ip{\vx}{\vy} =
        \vy^{\top} \vx\).
    \item El producto interno Frobenius de matrices es \(\ip{\vX}{\vY}
        = \tr(\vY^{\top}\vX)\).
    \item Un vector o matriz de elementos unitarios se denota \(\vone\) 
        (la dimensión debería ser obvia por contexto). De manera similar, 
        un vector o matriz de elementos nulos se denota \(\vzero\). La matriz 
        identidad se denota como \(\vI\).
\end{itemize}

\section*{Notación de probabilidades}

\begin{itemize}
    \item La medida de probabilidad base es \(\Pr\), es decir, la
        probabilidad de que ocurra un evento \(A\) es \(\Pr[A]\). Podemos
        especificar la distribución de una variable aleatoria usando
        subíndices, es decir, \(\Pr_{\vx \sim \mu}[\vx \in S]\), describe
        las probabilidades asociadas con una variable aleatoria $\vx$ con
        distribución (medida) $\mu$. Si dos variables aleatorias $\vx$ y $\vx'$
        tienen la misma distribución, escribimos $\vx \equid \vx'$.
    \item El operador esperanza es \(\Ex\), con la misma advertencia sobre el 
        uso de subíndices. 
    \item El operador covarianza es \(\Cov\), con la misma advertencia sobre 
        el uso de subíndices.
    \item El operador correlación es \(\Corr\), con la misma advertencia sobre 
        el uso de subíndices.
    \item  en ciertos contextos de modelamiento probabilístico donde 
        se supone una densidad (especialmente en
        \Cref{ch:denoising,ch:conditional-inference}), también usamos la 
        notación $p_{\vx}$ para la densidad de una variable aleatoria $\vx$. 
        Esto se extiende a las densidades condicionales, p.ej. 
        $p_{\vx \mid \vy}$ para la densidad de $\vx$ condicionada por $\vy$.
    \item En general usamos letras griegas para denotar las realizaciones
        de variables aleatorias (en contraste on las propias variables 
        aleatorias). Por ejemplo, $p_{\vx}(\vxi)$, $p_{\vx \mid
        \vy}(\vxi \mid \vnu)$, $\Ex[\vx \mid \vy=\vnu]$, etc.
    \item El conjunto de distribuciones probabilísticas sobre un conjunto finito        %
        %\footnote{Esto
        %    debería ser un espacio medible para evitar problemas en teoría
        %    de la meddida y teoría de conjuntos. Si usted no ha estudiado
        %    teoría de la medida puede ignorar esto con relativa seguridad.}
        \(\cX\) se denota
        \(\Delta(\cX)\). (Para \(\cX = [n]\), esto es el simplex probabilistico
        que se puede incrustar en \(\R^{n}\)).
    \item La distribución gaussiana con media \(\vmu\) y covarianza
        \(\vSigma\) se denota \(\dNorm(\vmu, \vSigma)\).
    \item La distribución uniforme sobre un conjunto compacto \(\cX\) se
        denota \(\dUnif(\cX)\).
\end{itemize}



\section*{Notación de aprendizaje automático}

\begin{itemize}
    \item Los codificadores y reductores de ruido normalmente se denotan por 
        \(f\) y usualmente tienen parámetros \(\theta \in \Theta\). Usualmente
        escribiremos las características como \(\vz = f_{\theta}(\vx)\).
    \item Los decodificadores normalmente se denotan por \(g\) y usualmente 
        tienen parámetros \(\eta\). Usualmente escribiremos la auto-codificación
        como \(\hat{\vx} = g_{\eta}(\vz)\) correspondiente a \(\vx\).
    \item Cuando \(f\) y \(g\) se implementan como redes neurales,
        tendrán el mismo número de capas sin pérdida de generalidad;
        we write them as \(f = f^{L} \circ f^{L - 1}
        \circ \cdots \circ f^{2} \circ f^{1} \circ f^{\pre}\) and \(g
            = g^{\post} \circ g^{L} \circ g^{ L - 1} \circ \cdots \circ
        g^{2} \circ g^{1}\).
    \item We write the features at the input to layer \(\ell\) as
        \(\vz^{\ell}\) such that \(\vz^{\ell + 1} =
        f^{\ell}(\vz^{\ell})\) with \(\vz^{1} = f^{\pre}(\vx)\) and
        \(\vz = \vz^{L + 1}\). Similarly the autoencoding features
        are \(\hat{\vx}^{\ell}\) with \(\hat{\vx}^{\ell + 1} =
        g^{\ell}(\hat{\vx}^{\ell})\) with \(\hat{\vx}^{1} = \vz\) and
        \(\hat{\vx} = g^{\post}(\hat{\vx}^{L + 1})\).
    \item For denoising, optimization, and processes which have a
        continuous time index, the time is almost always a subscript,
        e.g., \(\vx_{t}\). Discrete sequences can have the index be a
        superscript or subscript.
\end{itemize}

\section*{Modeling and Optimization Notation}
\begin{itemize}
    \item As in the Probability and Machine Learning Notation sections above,
        when we have a model, say for the realizations of
        a data distribution supported on $\bR^D$, we always denote it with
        ``plain'' accented variables, say $\vx$.
    \item In cases where we assume our data $\vx$ is generated by a model from
        a restricted class of \textit{parametric models}, say with
        a parameter $\vmu$, we also denote the \textit{true parameter} with
        plain accenting.\footnote{This could arise either due to us making an
        assumption about our data, or due to our data being synthetically
        generated. This convention allows us to adopt a concise notation across
        both `probabilistic' and `analytical' modeling setups. Readers more
        familiar with the latter kinds of frameworks may be used to using
        notation such as $\vmu_{\sharp}$ or $\vmu_o$ for this case; contrast
        with our accented decision variable notation.}
    \item \textit{Decision variables} (or learnable parameters) of a model fit
        to observed data are denoted with ``tilde'' accenting, e.g.\
        $\tilde{\vx}$, in cases where there is an underlying model $\vx$, or
        with plain notation in cases where no confusion is possible (e.g., if
        an empirically-fit model for \textit{nonparametric} data $\vx$ involves
        ``mean'' parameters $\vmu$).
    \item \textit{Optimal solutions} to optimization problems are denoted with
        ``star'' accenting, either as a superscript or a subscript (say
        $\vx_\star$ or $\vx^\star$).
    \item \textit{Estimators} for data that correspond to a specific statistical
        model, especially for the mean of that statistical model, are denoted
        with ``bar'' accenting, say $\bar{\vx}$ (especially for minimum
        mean-squared error denoising in \Cref{ch:denoising}).
    \item \textit{Approximations} associated to computational procedures for
        modeling data are denoted with ``hat'' accenting, say $\hat{\vx}$ for
        the output of an autoencoder or a generative model that approximates
        data $\vx$.\footnote{We generally distinguish these cases from optimal
        solutions $\vx_{\star}$, although in some special cases they are
        equal (e.g., in several examples in \Cref{ch:classic-models}). Think of
        the distinction between the parameter $\vmu$ of a model, which might be
        fit with optimization, and the results of sampling with that learned
        parameter.}
\end{itemize}
